%----------------------------------------------------------------------------------------
%	PACKAGES AND OTHER DOCUMENT CONFIGURATIONS
%----------------------------------------------------------------------------------------

\documentclass[12pt,a4paper,notitlepage,twoside,openright]{report}

\usepackage{xgreek} % Required for greek hyphenation
\usepackage[usenames,dvipsnames]{color} % Required for custom colors
%\usepackage{fontspec} % Required for fonts configuration
\usepackage{fancyhdr} % Required for custom headers
\usepackage{lastpage} % Required to determine the last page for the footer
%\usepackage{extramarks} % Required for headers and footers
\usepackage{graphicx} % Required to insert images
\usepackage[margin=4pt,font=normalsize,labelfont=bf,textfont=it]{caption} % captions
\usepackage{subcaption} % subfigures
\usepackage{listings} % Required for insertion of code
\usepackage{unicode-math} % Required to display math
\usepackage{tocloft} % Used to display dots in the table of contents
\usepackage{etoolbox} % Used to customize fancyhdr macros to colorize rule(s)
\usepackage{setspace}
\usepackage[colorlinks=true, urlcolor=LinksGray]{hyperref} %hyperlinks
\usepackage[notlof,notlot,nottoc]{tocbibind} %bibliography is not a numbered chapter
\usepackage[top=1.25in, bottom=1in, left=1in, right=1in]{geometry}
\usepackage{array} % for use with tables
\usepackage{mdwlist} % for compact lists
\usepackage{metalogo} % XeLaTeX logo
% Fonts
\usepackage{libertine} %Font -> Linux Libertine
\usepackage{courier} %Font -> Courier 
%\defaultfontfeatures{SizeFeatures={Size=12},Scale=MatchLowercase,Script=Greek}
%\setmainfont[Mapping=tex-text]{Linux Libertine O}
%\setsansfont[Mapping=tex-text]{Liberaton Sans}
\setmonofont[SizeFeatures={Size=10}]{Courier}
% \setmathfont[Scale=MatchUppercase]{Asana Math}

%----------------------------------------------------------------------------------------
%	TABLE OF CONTENTS
%----------------------------------------------------------------------------------------

\hypersetup{linkcolor=blue,citecolor=blue,linktoc=page} %hyperlink options

%----------------------------------------------------------------------------------------
%	PAGE STYLE
%----------------------------------------------------------------------------------------

\pagestyle{fancy}
\lhead[\rm \bf \thepage]{\fancyplain{}{\sl{\rightmark}}}
\chead{}
\rhead{\rm \bf \thepage}{\fancyplain{}{\sl{\leftmark }}}
\rfoot{}
\cfoot{}
\lfoot{}
%\renewcommand{\headrulewidth}{0.4pt}
\renewcommand{\footrulewidth}{0.4pt}

%----------------------------------------------------------------------------------------
%	CODE INCLUSION CONFIGURATION
%----------------------------------------------------------------------------------------

\definecolor{CommentsDarkGreen}{rgb}{0.0,0.4,0.0} % This is the color used for comments
\definecolor{NumbersLightGray}{rgb}{0.6,0.6,0.6} % This is the color used for numbers
\lstloadlanguages{Java} % Load Java syntax for listings, for a list of other languages supported see: ftp://ftp.tex.ac.uk/tex-archive/macros/latex/contrib/listings/listings.pdf
\lstset{% Structure
        frame=bottomline, % Single frame around code
        captionpos=b, % Place caption to the bottom
        breaklines=true, % Wrap if line is too big
        postbreak=\raisebox{0ex}[0ex][0ex]{\ensuremath{\rightdowncurvedarrow\space}},
       % Numbering
        numbers=left, % Line numbers on left
        firstnumber=1, % Line numbers start with line 1
        numberstyle=\footnotesize\color{NumbersLightGray}, % Line numbers are light gray and small
        numberfirstline=true, % Always number the first line
        stepnumber=1, % Line numbers go in steps of 1
        % Code
        language=Java, % Use Java in this example
        basicstyle=\tiny\ttfamily, % Use small true type font
        keywordstyle=\color{Blue}\bf\ttfamily, % keywords bold and blue
        identifierstyle=, % Nothing special about identifiers
        commentstyle=\usefont{T1}{pcr}{m}{sl}\color{CommentsDarkGreen}\small, % Code comments small dark green courier font
        stringstyle=\color{Purple}, % Strings are purple
        showstringspaces=false, % Don't put marks in string spaces
        tabsize=4, % 4 spaces per tab
        keepspaces=true, % keep spaces in code, usefull for indentation
        morecomment=[l][\color{Blue}]{...} % Line continuation (...) like blue comment
}

% Creates a new command to include a snippet from file
%   the first parameter is for the label and caption of the listing
%   the second parameter is the startline of the snippet
%   the third parameter is the endtline of the snippet
%   the fourth parameter is the name of the file (without .java)
\newcommand{\jsnippet}[4]{
\lstinputlisting[language=Java, firstline=#1, lastline=#2, #3]{#4}
}

\newcommand{\code}[1]{
\ttfamily{#1}\rmfamily\
}

\newcommand{\HRule}{\rule{\linewidth}{0.5mm}} % New command to make the lines in the title page


%----------------------------------------------------------------------------------------
%	COLORS
%----------------------------------------------------------------------------------------

\definecolor{LightGray}{rgb}{0.45,0.45,0.45}
\definecolor{LinksGray}{rgb}{0.3,0.3,0.3}

%----------------------------------------------------------------------------------------
%	NAME AND CLASS SECTION
%----------------------------------------------------------------------------------------

\newcommand{\hmwkTitle}{Ανάπτυξη Εφαρμογής Android για την Εύρεση και Απεικόνιση Γεωγραφικών Δεδομένων\\με συγχρονές τεχνολογίες web} % Assignment title
\newcommand{\hmwkType}{Πτυχιακή Εργασία}
\newcommand{\hmwkAuthorName}{Ορέστης\ Μ.\ Κτενάς\ [Π03078]} % Author name 
\newcommand{\hmwkAuthorEmail}{orestiskt@gmail.com} % Author email
\newcommand{\hmwkSupervisor}{Μαρία Βίρβου} % Supervisor

\author{\normalsize\textbf{\hmwkAuthorName}\\ \normalsize\textcolor{LightGray}{\hmwkAuthorEmail}}
\date{\normalsize{Πειραιάς\ 2014}} % Insert date here if you want it to appear below author's name

%----------------------------------------------------------------------------------------

\begin{document}
\pagenumbering{roman}
\begin{titlepage}
\begin{minipage}{0.17\textwidth}
	\includegraphics[width=2cm]{images/logo-unipi.png}
\end{minipage}
\begin{minipage}{0.7\textwidth}
	\begin{flushleft}
		\hspace*{0.5cm}\large ΠΑΝΕΠΙΣΤΗΜΙΟ ΠΕΙΡΑΙΩΣ \normalsize \\
		\hspace*{0.5cm}ΤΜΗΜΑ ΠΛΗΡΟΦΟΡΙΚΗΣ \\
		\hspace*{0.5cm}ΚΑΤΕΥΘΥΝΣΗ ΤΕΧΝΟΛΟΓΙΑΣ ΛΟΓΙΣΜΙΚΟΥ\\
		\hspace*{0.5cm}ΚΑΙ ΕΥΦΥΩΝ ΣΥΣΤΗΜΑΤΩΝ
	\end{flushleft}
\end{minipage}
\begin{center}
	\vspace{1.2in}
	\noindent\rule{15cm}{1.5pt} \\
	\bigskip
	\LARGE{\textbf{\hmwkTitle}}\\
	\noindent\rule{15cm}{1.5pt} \\
	\vspace{0.5in}
	\normalsize{\hmwkType}\\
	\vspace{0.5in}
	\normalsize\textbf{\hmwkAuthorName}\\ \normalsize\textcolor{LightGray}{\hmwkAuthorEmail}\\
	\vspace{3in}
	\normalsize{Επιβλέπων:\ {\hmwkSupervisor}}\\
	\vspace{0.5in}
	\normalsize{Πειραιάς 2015}	
\end{center}

\end{titlepage}

\thispagestyle{empty}

\newpage
\vspace*{2.0in}
\begin{center}
\textit{\large Ευχαριστίες} 
\end{center}
\bigskip
{\it
Αρχικά θα ήθελα να ευχαριστήσω την κ. Μ. Βίρβου για την ευκαιρία που μου έδωσε να ασχοληθώ με ένα τόσο ενδιαφέρον και σύγχρονο θέμα όπως είναι η ανάπτυξη εφαρμογών Android. Επιπλέον θα ήθελα να ευχαριστήσω θερμά τον διδάκτορα κ. Ευθύμιο Αλέπη για το χρόνο που μου διέθεσε, την προθυμία και την καθοδήγησή του κατά την εκπόνηση της εργασίας. Θα ήθελα επίσης να ευχαριστήσω τους φίλους και συνεργάτες Α. Καραλιά και Α. Ψεγιαννάκη για τις πολύτιμες συμβουλές και τη βοήθειά τους στα τεχνικά θέματα της εργασίας. Τέλος θα ήθελα να ευχαριστήσω την οικογένειά μου και ιδιαίτερα τους γονείς μου για τη ηθική υποστήριξη και πίστη καθ'όλη τη διάρκεια των σπουδών μου.
}

\newpage
\tableofcontents
\clearpage \pagenumbering{arabic} \pagestyle{fancy}
\chapter{Εισαγωγή}
\lhead{Κεφάλαιο 1. \emph{Εισαγωγή}}

\begin{section}{Context-aware εφαρμογές}
\bigskip
Ένα σύστημα θεωρείται context-aware αν χρησιμοποιεί τα συμφραζόμενα (context) εκτέλεσης για να παράσχει συσχετισμένη πληροφορία ή υπηρεσία στο χρήστη, όπου η σχετικότητα εξαρτάται από το στόχο του χρήστη. Τα συμφραζόμενα χωρίζονται σε τέσσερις ευρύτερες κατηγορίες:
\begin{itemize}
\item \emph{Υπολογιστικά συμφραζομενα} όπως συνδεσιμότητα σε δίκτυο, bandwidth επικοινωνίας, πόροι συστηματος. 
\item \emph{Συμφραζόμενα χρήστη} προφίλ του χρήστη, γεωγραφική θέση, δραστηριότητα, συναισθηματική κατάσταση.
\item \emph{Φυσικά συμφραζόμενα} φωτισμός, θόρυβος περιβάλλοντος, θερμοκρασία.
\item \emph{Χρονικό πλαίσιο} ώρα της ημέρας, εποχη του χρόνου. \cite{ctxt_aware_comp}
\end{itemize}
Οι σύγχρονες κινητές συσκευές, παρέχουν του υλικό και λογισμικό για να αντλήσουν πληροφορίες για τα συμφραζομενα και να υποστηρίξουν συσχέτιση περιεχομένου σε πραγματικό χρόνο. Αυτές οι συσκευές ενσωματώνουν τεχνολογίες ικανές να ανιχνεύσουν την κατάσταση, το περιβάλλον και την δραστηριότητα του χρήστη. Τεχνολογίες όπως το GPS, γυροσκόπιο, επιταχυνσιόμετρο, πυξίδα, αισθητήρες εγγύτητας, αισθητήρες φωτός. Το λογισμικό εκμεταλλευόμενο αυτές τις τεχνολογίες μπορεί να παρέχει στοχευμένες υπηρεσίες και να ενισχύσει την εμπειρία χρήσης.
\par
Η γεωγραφική θέση του χρήστη διαδραματίζει πρωταγωνιστικό ρόλο στην εξαγωγή συμφραζομένων μέσω της αναγνώρισης του περιβάλλοντος χρήσης λογισμικού. Με επίκεντρο τη θέση του χρήστη έχουν αναπτυχθεί υπηρεσίες αναζήτησης και απεικόνισης συσχετισμένου περιεχομένου. Η προβολή σημείων ενδιαφέροντος (Points of Interest ή POI) περιλαμβάνει στοιχεία του χάρτη που καταλαβάνουν σημειακές θέσεις σε αντίθεση με γραμμικά δεδομενα όπως δρόμοι και σύνοριακές γραμμές. Τέτοια σημεία χρησιμοποιούνται κατά κόρον σε εφαρμογές πλοήγησης. Τα σημεία ενδιαφέροντος λαμβάνουν διαφορετική σημασία, ανάλογα με το πεδίο εφαρμογής. Τέτοια γεωγραφικά δεδομένα μπορεί να περιλαμβάνουν δημόσια κτήρια, σταθμούς μέσων μαζικής μεταφοράς, σχολεία, νοσοκομεία, αρχαιολογικούς χώρους και πολλά άλλα.

\end{section}


\begin{section}{Στόχος της εργασίας}
\bigskip
Στόχος της εργασίας είναι να ερευνήσει και να υλοποιήσει ένα σύστημα που εκμεταλλεύεται τις σύγχρονες mobile τεχνολογίες εντοπισμού και συγκειμένου για να παρουσιάσει στο χρήστη συσχετισμένο περιεχόμενο. Για το σκοπό αυτό υλοποιήθηκε εφαρμογή για το λειτουργικό κινητών συσκευών Android. Σαν use case επιλέχθηκε η εύρεση πρατηριών καυσίμων. Η εφαρμογή δειγματοληπτεί τη γεωγραφική θέση του χρήστη για να τροφοδοτήσει με τα κοντινότερα πρατήρια. Παράλληλα υλοποιήθηκε υπηρεσία αποθήκευσης και παροχής δεδομένων μέσω του διαδικτύου βασισμένη στο Elasticsearch και το Spring framework.

\end{section}


\begin{section}{Διάρθρωση της εργασίας}
\bigskip
Αρχικά γίνεται μια εισαγωγή στο Android περιβάλλον. Στόχος αυτού του κεφαλαίου είναι να παρουσιάσει στον αναγνώστη, εν συντομία, το λειτουργικό, το περιβάλλον ανάπτυξης εφαρμογών και τις τεχνολογίες του API που χρησιμοποιήθηκαν. Περιγράφονται συνοπτικά τα επίπεδα του Android stack και γίνεται μια εισαγωγή στα βασικά μέρη μιας εφαρμογής Android. Τέλος αναλύονται οι υπηρεσίας τοποθεσίας.
\par
Στο τρίτο κεφάλαιο γίνεται μια συνοπτική παρουσίαση της εφαρμογής. Περιγράφεται η λειτουργικότητα και η διεπαφή και δίνεται μια σύντομη εικόνα της αρχιτεκτονικής.
\par
Στο τέταρτο κεφάλαιο γίνεται η τεχνική ανάλυση της υλοποίησης της εφαρμογής. Παρουσιάζονται διεξοδικά οι πτυχές της υλοποίησης με αναφορές στις προγραμματιστικές λεπτομέρειες.
\par
Στο πέμπτο κεφάλαιο αναλύεται το κομμάτι του server. Πρώτα γίνεται μια εισαγωγή στο Elasticsearch και την ορολογια που το συνοδεύει. Ύστερα γίνεται η τεχνική ανάλυση του Indexer και του client.
\par
Τέλος γίνεται η αποτίμηση της εργασίας και δίνονται προτάσεις βελτίωσης και επέκτασης της εφαρμογής.

\end{section}


\begin{section}{Εργαλεία που χρησιμοποιήθηκαν}
\bigskip
Για την ανάπτυξη της Android εφαρμογής χρησιμοποιήθηκαν τα εξής εργαλεία:
\begin{itemize}
\item Android SDK Platform έκδοση 19.
\item η βιβλιοθήκη Google Play Services για τη χρήση των Google Maps κα Location Services.
\item η βιβλιοθήκη Spring for Android για τα REST calls στον απομακρυσμένο server. \cite{spring_and_proj}
\item η βιβλιοθήκη Jackson για marshalling/unmarshalling των JSON documents. \cite{jackson}
\item η βιβλιοθήκη Google Maps Android API utility library για τη σχεδίαση των σημείων. \cite{map_utils}
\item η ανάπτυξη έγινε στο Android Studio 0.8.9 IDE με το εργαλείο Gradle για dependency management.
\end{itemize}

Στο back-end χρησιμοποιήθηκαν τα εξής:
\begin{itemize}
\item Spring Boot με embedded Tomcat server. \cite{spring_boot_proj}
\item Elasticsearch server (1.3.2) ως NoSQL data-store και search server. \cite{es_home}
\item το Java API του elasticsearch (1.3.2) για indexing και search. \cite{es_java}
\item η ανάπτυξη έγινε στο Netbeans 8.0.1 IDE με το εργαλείο Maven για dependency management.
\end{itemize}

Τα δεδομένα των πρατηρίων μεταφορτώθηκαν από τον κρατικό ιστότοπο \url{data.gov.gr} σε μορφή Microsoft Excel (.xlsx) και αφορούν το μήνα Σεπτέμβριο 2013.

Το παρόν στοιχειοθετήθηκε με \XeLaTeX.

\end{section}

\begin{thebibliography}{99}

\bibitem{ctxt_aware_comp} Schlit B, Adams N, Want R, 1994, Context-aware computing applications, proceedings of the Mobile Computing Systems and Applications (WMCSA) conference, [1994]

\bibitem{oha_open_platform} Open Handset Alliance, 2007, Industry Leaders Announce Open Platform for Mobile Devices, Available from: <\url{http://www.openhandsetalliance.com/press_110507.html}>, [5 November 2007]

\bibitem{oha_sdk_release} Open Handset Alliance, 2007, Open Handset Alliance Releases Android SDK, Available from: <\url{http://www.openhandsetalliance.com/press_111207.html}>, [12 November 2007]

\bibitem{jc_kernel} Corbet J, 2011, Bringing Android closer to the mainline,  Available from: <\url{https://lwn.net/Articles/472984/}>, [20 December 2011]

\bibitem{kernel_newbies} Kernel Newbies Wiki, 2012, Linux 3.3, Tech. Report, Available from: <\url{http://kernelnewbies.org/Linux_3.3}>, [March 2012]

\bibitem{and_dalvik} Android Technical Information, Dalvik Technical Information, Available from: <\url{https://source.android.com/devices/tech/dalvik/index.html}>

\bibitem{dalvik_serg} Khan So, Khan Sh, Banuri H, Nauman M, Alam M, 2009, The Dalvik Virtual Machine Architecture, Security Engineering Group, Institute of Management Sciences Peshawar, Pakistan, Tech. Report, [November 2009]

\bibitem{dalvik_ehringer} Ehringer D, 2010, The Dalvik Virtual Machine Architecture, Technical Report, [March 2010]

\bibitem{and_art} Android Techincal Information, Introducing ART, Available from: <\url{https://source.android.com/devices/tech/dalvik/art.html}>

\bibitem{art_fru} Frumusanu A, 2014, A Closer Look at Android RunTime (ART) in Android L, Available from: <\url{http://www.anandtech.com/show/8231/a-closer-look-at-android-runtime-art-in-android-l}>, [1 July 2014]

\bibitem{apache_harmony} Apache Harmony, Apache Harmony home page, <\url{http://harmony.apache.org/}>

\bibitem{and_fund} Android API Guide, Application Fundamentals, Available from: <\url{http://developer.android.com/guide/components/fundamentals.html}>

\bibitem{and_act_life} Android API Guide, Managing the Activity Lifecycle, Available from: <\url{http://developer.android.com/training/basics/activity-lifecycle/starting.html}>

\bibitem{and_frag} Android API Guide, Fragments. Available from: <\url{http://developer.android.com/guide/components/fragments.html}>

\bibitem{and_cont} Android API Guide, Content Providers, Available from: <\url{http://developer.android.com/guide/topics/providers/content-providers.html}>ή

\bibitem{and_maps} Google Maps Android API v2 , Available from: <\url{https://developers.google.com/maps/documentation/android/}>

\bibitem{and_maps} Location APIs, Available from: <\url{https://developer.android.com/google/play-services/location.html}>

\bibitem{and_sdk}Android SDK, Available from: <\url{http://developer.android.com/sdk/index.html}>

\bibitem{and_guide} Android API Guide, Available from: <\url{http://developer.android.com/guide/index.html}>

\bibitem{and_train}Android Training, Available from: <\url{http://developer.android.com/training/index.html}>

\bibitem{and_ref} Android API reference, Available from: <\url{http://developer.android.com/reference/packages.html}>

\bibitem{murphy_busy_coder} Murphy M, 2012, The busy coder's guide to Android development, v4.2 [U.S.]: CommonsWare, LLC.

\bibitem{cookbook} Darwin I, 2012, Android cookbook, 1st ed, Sebastopol, CA: O'Reilly

\bibitem{gor_thread} Göransson A, 2014, Efficient Android Threading:  Asynchronous Processing Techniques for Android Applications, 1st ed. Sebastopol, CA: O'Reilly Media.

\bibitem{and_thread} Android API Guide, Processes and Threads, Available from: <\url{http://developer.android.com/guide/components/processes-and-threads.html}>

\bibitem{and_arch_serg} Khan Sh, Banuri H, Nauman M, Khan So, Alam M, 2010, Analysis report on Android Application Framework and existing Security Architecture', Security Engineering Group, Institute of Management Sciences Peshawar, Pakistan, Tech. Report[February 2010]

\bibitem{map_utils} Google Maps Android API utility library, <\url{http://googlemaps.github.io/android-maps-utils/}>

\bibitem{spring_and_proj} Spring for Android Project Page, <\url{http://projects.spring.io/spring-android/}>

\bibitem{jackson}Jackson Project Home, <\url{https://github.com/FasterXML/jackson}>

\bibitem{spring_and} Spring for Android Reference Documentation, Available from: <\url{http://docs.spring.io/spring-android/docs/current/reference/htmlsingle/}>

\bibitem{ap_lucene} Apache Lucene Project Page, <\url{http://lucene.apache.org/}>

\bibitem{es_home} Elasticsearch home page, <\url{http://www.elasticsearch.org/}>

\bibitem{es_ref} Elastisearch Reference Guide, Available from: <\url{http://www.elasticsearch.org/guide/en/elasticsearch/reference/current/index.html}>

\bibitem{es_java} Elastisearch Java API Reference Guide, Available from:
<\url{http://www.elasticsearch.org/guide/en/elasticsearch/client/java-api/current/}>

\bibitem{es_inaction} Gheorghe R, Hinman M L, 2014, Elasticsearch In Action, Manning Publications Co. [2014]

\bibitem{geocoding} The Google Geocoding API Guide, Available from: <\url{https://developers.google.com/maps/documentation/geocoding/}>

\bibitem{spring_boot_proj} Spring Boot Project Page, <\url{http://projects.spring.io/spring-boot/}>

\bibitem{spring_boot} Spring Boot Reference Guide, Available from: <\url{http://docs.spring.io/spring-boot/docs/1.1.6.RELEASE/reference/htmlsingle/}>

\end{thebibliography}
\end{document}
